\documentclass[11pt,twoside,final,openright]{report}
\usepackage{a4,graphicx,html,parskip,setspace,times,xspace,url}
\setstretch{1.15}

\renewcommand{\ttdefault}{pcr}

\def\Xend{{Xend}\xspace}
\def\xend{{xend}\xspace}

\latexhtml{\renewcommand{\path}[1]{{\small {\tt #1}}}}{\renewcommand{\path}[1]{{\tt #1}}}


\begin{document}

% TITLE PAGE
\pagestyle{empty}
\begin{center}
\vspace*{\fill}
\includegraphics{figs/xenlogo.eps}
\vfill
\vfill
\vfill
\begin{tabular}{l}
{\Huge \bf Users' Manual} \\[4mm]
{\huge Xen v3.0} \\[80mm]
\end{tabular}
\end{center}

{\bf DISCLAIMER: This documentation is always under active development
and as such there may be mistakes and omissions --- watch out for
these and please report any you find to the developers' mailing list,
xen-devel@lists.xensource.com. The latest version is always available
on-line. Contributions of material, suggestions and corrections are
welcome.}

\vfill
\clearpage


% COPYRIGHT NOTICE
\pagestyle{empty}

\vspace*{\fill}

Xen is Copyright \copyright  2002-2005, University of Cambridge, UK, XenSource
Inc., IBM Corp., Hewlett-Packard Co., Intel Corp., AMD Inc., and others.  All
rights reserved.

Xen is an open-source project.  Most portions of Xen are licensed for copying
under the terms of the GNU General Public License, version 2.  Other portions
are licensed under the terms of the GNU Lesser General Public License, the
Zope Public License 2.0, or under ``BSD-style'' licenses.  Please refer to the
COPYING file for details.

\cleardoublepage


% TABLE OF CONTENTS
\pagestyle{plain}
\pagenumbering{roman}
{ \parskip 0pt plus 1pt
  \tableofcontents }
\cleardoublepage


% PREPARE FOR MAIN TEXT
\pagenumbering{arabic}
\raggedbottom
\widowpenalty=10000
\clubpenalty=10000
\parindent=0pt
\parskip=5pt
\renewcommand{\topfraction}{.8}
\renewcommand{\bottomfraction}{.8}
\renewcommand{\textfraction}{.2}
\renewcommand{\floatpagefraction}{.8}
\setstretch{1.1}


%% Chapter Introduction moved to introduction.tex
\chapter{Introduction}


Xen is an open-source \emph{para-virtualizing} virtual machine monitor
(VMM), or ``hypervisor'', for the x86 processor architecture. Xen can
securely execute multiple virtual machines on a single physical system
with close-to-native performance.  Xen facilitates enterprise-grade
functionality, including:

\begin{itemize}
\item Virtual machines with performance close to native hardware.
\item Live migration of running virtual machines between physical hosts.
\item Up to 32 virtual CPUs per guest virtual machine, with VCPU hotplug.
\item x86/32, x86/32 with PAE, and x86/64 platform support.
\item Intel Virtualization Technology (VT-x) for unmodified guest operating systems (including Microsoft Windows).
\item Excellent hardware support (supports almost all Linux device
  drivers). 
\end{itemize}


\section{Usage Scenarios}

Usage scenarios for Xen include:

\begin{description}
\item [Server Consolidation.] Move multiple servers onto a single
  physical host with performance and fault isolation provided at the
  virtual machine boundaries.
\item [Hardware Independence.] Allow legacy applications and operating 
  systems to exploit new hardware.
\item [Multiple OS configurations.] Run multiple operating systems
  simultaneously, for development or testing purposes.
\item [Kernel Development.] Test and debug kernel modifications in a
  sand-boxed virtual machine --- no need for a separate test machine.
\item [Cluster Computing.] Management at VM granularity provides more
  flexibility than separately managing each physical host, but better
  control and isolation than single-system image solutions,
  particularly by using live migration for load balancing.
\item [Hardware support for custom OSes.] Allow development of new
  OSes while benefiting from the wide-ranging hardware support of
  existing OSes such as Linux.
\end{description}


\section{Operating System Support}

Para-virtualization permits very high performance virtualization, even
on architectures like x86 that are traditionally very hard to
virtualize.

This approach requires operating systems to be \emph{ported} to run on
Xen. Porting an OS to run on Xen is similar to supporting a new
hardware platform, however the process is simplified because the
para-virtual machine architecture is very similar to the underlying
native hardware. Even though operating system kernels must explicitly
support Xen, a key feature is that user space applications and
libraries \emph{do not} require modification.

With hardware CPU virtualization as provided by Intel VT and AMD
SVM technology, the ability to run an unmodified guest OS kernel
is available.  No porting of the OS is required, although some
additional driver support is necessary within Xen itself.  Unlike
traditional full virtualization hypervisors, which suffer a tremendous
performance overhead, the combination of Xen and VT or Xen and
Pacifica technology complement one another to offer superb performance
for para-virtualized guest operating systems and full support for
unmodified guests running natively on the processor.  Full support for
VT and Pacifica chipsets will appear in early 2006.

Paravirtualized Xen support is available for increasingly many
operating systems: currently, mature Linux support is available and
included in the standard distribution.  Other OS ports---including
NetBSD, FreeBSD and Solaris x86 v10---are nearing completion. 


\section{Hardware Support}

Xen currently runs on the x86 architecture, requiring a ``P6'' or
newer processor (e.g.\ Pentium Pro, Celeron, Pentium~II, Pentium~III,
Pentium~IV, Xeon, AMD~Athlon, AMD~Duron). Multiprocessor machines are
supported, and there is support for HyperThreading (SMT).  In 
addition, ports to IA64 and Power architectures are in progress.

The default 32-bit Xen supports up to 4GB of memory. However Xen 3.0
adds support for Intel's Physical Addressing Extensions (PAE), which
enable x86/32 machines to address up to 64 GB of physical memory.  Xen
3.0 also supports x86/64 platforms such as Intel EM64T and AMD Opteron
which can currently address up to 1TB of physical memory.

Xen offloads most of the hardware support issues to the guest OS
running in the \emph{Domain~0} management virtual machine. Xen itself
contains only the code required to detect and start secondary
processors, set up interrupt routing, and perform PCI bus
enumeration. Device drivers run within a privileged guest OS rather
than within Xen itself. This approach provides compatibility with the
majority of device hardware supported by Linux. The default XenLinux
build contains support for most server-class network and disk
hardware, but you can add support for other hardware by configuring
your XenLinux kernel in the normal way.


\section{Structure of a Xen-Based System}

A Xen system has multiple layers, the lowest and most privileged of
which is Xen itself.

Xen may host multiple \emph{guest} operating systems, each of which is
executed within a secure virtual machine. In Xen terminology, a
\emph{domain}. Domains are scheduled by Xen to make effective use of the
available physical CPUs. Each guest OS manages its own applications.
This management includes the responsibility of scheduling each
application within the time allotted to the VM by Xen.

The first domain, \emph{domain~0}, is created automatically when the
system boots and has special management privileges. Domain~0 builds
other domains and manages their virtual devices. It also performs
administrative tasks such as suspending, resuming and migrating other
virtual machines.

Within domain~0, a process called \emph{xend} runs to manage the system.
\Xend\ is responsible for managing virtual machines and providing access
to their consoles. Commands are issued to \xend\ over an HTTP interface,
via a command-line tool.


\section{History}

Xen was originally developed by the Systems Research Group at the
University of Cambridge Computer Laboratory as part of the XenoServers
project, funded by the UK-EPSRC\@.

XenoServers aim to provide a ``public infrastructure for global
distributed computing''. Xen plays a key part in that, allowing one to
efficiently partition a single machine to enable multiple independent
clients to run their operating systems and applications in an
environment. This environment provides protection, resource isolation
and accounting. The project web page contains further information along
with pointers to papers and technical reports:
\path{http://www.cl.cam.ac.uk/xeno}

Xen has grown into a fully-fledged project in its own right, enabling us
to investigate interesting research issues regarding the best techniques
for virtualizing resources such as the CPU, memory, disk and network.
Project contributors now include XenSource, Intel, IBM, HP, AMD, Novell,
RedHat.

Xen was first described in a paper presented at SOSP in
2003\footnote{\tt
  http://www.cl.cam.ac.uk/netos/papers/2003-xensosp.pdf}, and the first
public release (1.0) was made that October. Since then, Xen has
significantly matured and is now used in production scenarios on many
sites.

\section{What's New}

Xen 3.0.0 offers:

\begin{itemize}
\item Support for up to 32-way SMP guest operating systems
\item Intel (Physical Addressing Extensions) PAE to support 32-bit
  servers with more than 4GB physical memory
\item x86/64 support (Intel EM64T, AMD Opteron)
\item Intel VT-x support to enable the running of unmodified guest
operating systems (Windows XP/2003, Legacy Linux)
\item Enhanced control tools
\item Improved ACPI support
\item AGP/DRM graphics
\end{itemize}


Xen 3.0 features greatly enhanced hardware support, configuration
flexibility, usability and a larger complement of supported operating
systems.  This latest release takes Xen a step closer to being the 
definitive open source solution for virtualization.



\part{Installation}

%% Chapter Basic Installation
\chapter{Basic Installation}

The Xen distribution includes three main components: Xen itself, ports
of Linux and NetBSD to run on Xen, and the userspace tools required to
manage a Xen-based system. This chapter describes how to install the
Xen~3.0 distribution from source. Alternatively, there may be pre-built
packages available as part of your operating system distribution.


\section{Prerequisites}
\label{sec:prerequisites}

The following is a full list of prerequisites. Items marked `$\dag$' are
required by the \xend\ control tools, and hence required if you want to
run more than one virtual machine; items marked `$*$' are only required
if you wish to build from source.
\begin{itemize}
\item A working Linux distribution using the GRUB bootloader and running
  on a P6-class or newer CPU\@.
\item [$\dag$] The \path{iproute2} package.
\item [$\dag$] The Linux bridge-utils\footnote{Available from {\tt
      http://bridge.sourceforge.net}} (e.g., \path{/sbin/brctl})
\item [$\dag$] The Linux hotplug system\footnote{Available from {\tt
      http://linux-hotplug.sourceforge.net/}} (e.g.,
      \path{/sbin/hotplug} and related scripts).  On newer distributions,
      this is included alongside the Linux udev system\footnote{See {\tt
      http://www.kernel.org/pub/linux/utils/kernel/hotplug/udev.html/}}.
\item [$*$] Build tools (gcc v3.2.x or v3.3.x, binutils, GNU make).
\item [$*$] Development installation of zlib (e.g.,\ zlib-dev).
\item [$*$] Development installation of Python v2.2 or later (e.g.,\
  python-dev).
\item [$*$] \LaTeX\ and transfig are required to build the
  documentation.
\end{itemize}

Once you have satisfied these prerequisites, you can now install either
a binary or source distribution of Xen.

\section{Installing from Binary Tarball}

Pre-built tarballs are available for download from the XenSource downloads
page:
\begin{quote} {\tt http://www.xensource.com/downloads/}
\end{quote}

Once you've downloaded the tarball, simply unpack and install:
\begin{verbatim}
# tar zxvf xen-3.0-install.tgz
# cd xen-3.0-install
# sh ./install.sh
\end{verbatim}

Once you've installed the binaries you need to configure your system as
described in Section~\ref{s:configure}.

\section{Installing from RPMs}
Pre-built RPMs are available for download from the XenSource downloads
page:
\begin{quote} {\tt http://www.xensource.com/downloads/}
\end{quote}

Once you've downloaded the RPMs, you typically install them via the 
RPM commands: 

\verb|# rpm -iv rpmname| 

See the instructions and the Release Notes for each RPM set referenced at:
  \begin{quote}
    {\tt http://www.xensource.com/downloads/}.
  \end{quote}
 
\section{Installing from Source}

This section describes how to obtain, build and install Xen from source.

\subsection{Obtaining the Source}

The Xen source tree is available as either a compressed source tarball
or as a clone of our master Mercurial repository.

\begin{description}
\item[Obtaining the Source Tarball]\mbox{} \\
  Stable versions and daily snapshots of the Xen source tree are
  available from the Xen download page:
  \begin{quote} {\tt \tt http://www.xensource.com/downloads/}
  \end{quote}
\item[Obtaining the source via Mercurial]\mbox{} \\
  The source tree may also be obtained via the public Mercurial
  repository at:
  \begin{quote}{\tt http://xenbits.xensource.com}
  \end{quote} See the instructions and the Getting Started Guide
  referenced at:
  \begin{quote}
    {\tt http://www.xensource.com/downloads/}
  \end{quote}
\end{description}

% \section{The distribution}
%
% The Xen source code repository is structured as follows:
%
% \begin{description}
% \item[\path{tools/}] Xen node controller daemon (Xend), command line
%   tools, control libraries
% \item[\path{xen/}] The Xen VMM.
% \item[\path{buildconfigs/}] Build configuration files
% \item[\path{linux-*-xen-sparse/}] Xen support for Linux.
% \item[\path{patches/}] Experimental patches for Linux.
% \item[\path{docs/}] Various documentation files for users and
%   developers.
% \item[\path{extras/}] Bonus extras.
% \end{description}

\subsection{Building from Source}

The top-level Xen Makefile includes a target ``world'' that will do the
following:

\begin{itemize}
\item Build Xen.
\item Build the control tools, including \xend.
\item Download (if necessary) and unpack the Linux 2.6 source code, and
  patch it for use with Xen.
\item Build a Linux kernel to use in domain~0 and a smaller unprivileged
  kernel, which can be used for unprivileged virtual machines.
\end{itemize}

After the build has completed you should have a top-level directory
called \path{dist/} in which all resulting targets will be placed. Of
particular interest are the two XenLinux kernel images, one with a
``-xen0'' extension which contains hardware device drivers and drivers
for Xen's virtual devices, and one with a ``-xenU'' extension that
just contains the virtual ones. These are found in
\path{dist/install/boot/} along with the image for Xen itself and the
configuration files used during the build.

%The NetBSD port can be built using:
%\begin{quote}
%\begin{verbatim}
%# make netbsd20
%\end{verbatim}\end{quote}
%NetBSD port is built using a snapshot of the netbsd-2-0 cvs branch.
%The snapshot is downloaded as part of the build process if it is not
%yet present in the \path{NETBSD\_SRC\_PATH} search path.  The build
%process also downloads a toolchain which includes all of the tools
%necessary to build the NetBSD kernel under Linux.

To customize the set of kernels built you need to edit the top-level
Makefile. Look for the line:
\begin{quote}
\begin{verbatim}
KERNELS ?= linux-2.6-xen0 linux-2.6-xenU
\end{verbatim}
\end{quote}

You can edit this line to include any set of operating system kernels
which have configurations in the top-level \path{buildconfigs/}
directory.

%% Inspect the Makefile if you want to see what goes on during a
%% build.  Building Xen and the tools is straightforward, but XenLinux
%% is more complicated.  The makefile needs a `pristine' Linux kernel
%% tree to which it will then add the Xen architecture files.  You can
%% tell the makefile the location of the appropriate Linux compressed
%% tar file by
%% setting the LINUX\_SRC environment variable, e.g. \\
%% \verb!# LINUX_SRC=/tmp/linux-2.6.11.tar.bz2 make world! \\ or by
%% placing the tar file somewhere in the search path of {\tt
%%   LINUX\_SRC\_PATH} which defaults to `{\tt .:..}'.  If the
%% makefile can't find a suitable kernel tar file it attempts to
%% download it from kernel.org (this won't work if you're behind a
%% firewall).

%% After untaring the pristine kernel tree, the makefile uses the {\tt
%%   mkbuildtree} script to add the Xen patches to the kernel.

%% \framebox{\parbox{5in}{
%%     {\bf Distro specific:} \\
%%     {\it Gentoo} --- if not using udev (most installations,
%%     currently), you'll need to enable devfs and devfs mount at boot
%%     time in the xen0 config.  }}

\subsection{Custom Kernels}

% If you have an SMP machine you may wish to give the {\tt '-j4'}
% argument to make to get a parallel build.

If you wish to build a customized XenLinux kernel (e.g.\ to support
additional devices or enable distribution-required features), you can
use the standard Linux configuration mechanisms, specifying that the
architecture being built for is \path{xen}, e.g:
\begin{quote}
\begin{verbatim}
# cd linux-2.6.12-xen0
# make ARCH=xen xconfig
# cd ..
# make
\end{verbatim}
\end{quote}

You can also copy an existing Linux configuration (\path{.config}) into
e.g.\ \path{linux-2.6.12-xen0} and execute:
\begin{quote}
\begin{verbatim}
# make ARCH=xen oldconfig
\end{verbatim}
\end{quote}

You may be prompted with some Xen-specific options. We advise accepting
the defaults for these options.

Note that the only difference between the two types of Linux kernels
that are built is the configuration file used for each. The ``U''
suffixed (unprivileged) versions don't contain any of the physical
hardware device drivers, leading to a 30\% reduction in size; hence you
may prefer these for your non-privileged domains. The ``0'' suffixed
privileged versions can be used to boot the system, as well as in driver
domains and unprivileged domains.

\subsection{Installing Generated Binaries}

The files produced by the build process are stored under the
\path{dist/install/} directory. To install them in their default
locations, do:
\begin{quote}
\begin{verbatim}
# make install
\end{verbatim}
\end{quote}

Alternatively, users with special installation requirements may wish to
install them manually by copying the files to their appropriate
destinations.

%% Files in \path{install/boot/} include:
%% \begin{itemize}
%% \item \path{install/boot/xen-3.0.gz} Link to the Xen 'kernel'
%% \item \path{install/boot/vmlinuz-2.6-xen0} Link to domain 0
%%   XenLinux kernel
%% \item \path{install/boot/vmlinuz-2.6-xenU} Link to unprivileged
%%   XenLinux kernel
%% \end{itemize}

The \path{dist/install/boot} directory will also contain the config
files used for building the XenLinux kernels, and also versions of Xen
and XenLinux kernels that contain debug symbols such as
(\path{xen-syms-3.0.0} and \path{vmlinux-syms-2.6.12.6-xen0}) which are
essential for interpreting crash dumps. Retain these files as the
developers may wish to see them if you post on the mailing list.


\section{Configuration}
\label{s:configure}

Once you have built and installed the Xen distribution, it is simple to
prepare the machine for booting and running Xen.

\subsection{GRUB Configuration}

An entry should be added to \path{grub.conf} (often found under
\path{/boot/} or \path{/boot/grub/}) to allow Xen / XenLinux to boot.
This file is sometimes called \path{menu.lst}, depending on your
distribution. The entry should look something like the following:

%% KMSelf Thu Dec  1 19:06:13 PST 2005 262144 is useful for RHEL/RH and
%% related Dom0s.
{\small
\begin{verbatim}
title Xen 3.0 / XenLinux 2.6
  kernel /boot/xen-3.0.gz dom0_mem=262144
  module /boot/vmlinuz-2.6-xen0 root=/dev/sda4 ro console=tty0
\end{verbatim}
}

The kernel line tells GRUB where to find Xen itself and what boot
parameters should be passed to it (in this case, setting the domain~0
memory allocation in kilobytes and the settings for the serial port).
For more details on the various Xen boot parameters see
Section~\ref{s:xboot}.

The module line of the configuration describes the location of the
XenLinux kernel that Xen should start and the parameters that should be
passed to it. These are standard Linux parameters, identifying the root
device and specifying it be initially mounted read only and instructing
that console output be sent to the screen. Some distributions such as
SuSE do not require the \path{ro} parameter.

%% \framebox{\parbox{5in}{
%%     {\bf Distro specific:} \\
%%     {\it SuSE} --- Omit the {\tt ro} option from the XenLinux
%%     kernel command line, since the partition won't be remounted rw
%%     during boot.  }}

To use an initrd, add another \path{module} line to the configuration,
like: {\small
\begin{verbatim}
  module /boot/my_initrd.gz
\end{verbatim}
}

%% KMSelf Thu Dec  1 19:05:30 PST 2005 Other configs as an appendix?

When installing a new kernel, it is recommended that you do not delete
existing menu options from \path{menu.lst}, as you may wish to boot your
old Linux kernel in future, particularly if you have problems.

\subsection{Serial Console (optional)}

Serial console access allows you to manage, monitor, and interact with
your system over a serial console.  This can allow access from another
nearby system via a null-modem (``LapLink'') cable or remotely via a serial
concentrator.

You system's BIOS, bootloader (GRUB), Xen, Linux, and login access must
each be individually configured for serial console access.  It is
\emph{not} strictly necessary to have each component fully functional,
but it can be quite useful.

For general information on serial console configuration under Linux,
refer to the ``Remote Serial Console HOWTO'' at The Linux Documentation
Project: \url{http://www.tldp.org} 

\subsubsection{Serial Console BIOS configuration}

Enabling system serial console output neither enables nor disables
serial capabilities in GRUB, Xen, or Linux, but may make remote
management of your system more convenient by displaying POST and other
boot messages over serial port and allowing remote BIOS configuration.

Refer to your hardware vendor's documentation for capabilities and
procedures to enable BIOS serial redirection.


\subsubsection{Serial Console GRUB configuration}

Enabling GRUB serial console output neither enables nor disables Xen or
Linux serial capabilities, but may made remote management of your system
more convenient by displaying GRUB prompts, menus, and actions over
serial port and allowing remote GRUB management.

Adding the following two lines to your GRUB configuration file,
typically either \path{/boot/grub/menu.lst} or \path{/boot/grub/grub.conf}
depending on your distro, will enable GRUB serial output.

\begin{quote} 
{\small \begin{verbatim}
  serial --unit=0 --speed=115200 --word=8 --parity=no --stop=1
  terminal --timeout=10 serial console
\end{verbatim}}
\end{quote}

Note that when both the serial port and the local monitor and keyboard
are enabled, the text ``\emph{Press any key to continue}'' will appear
at both.  Pressing a key on one device will cause GRUB to display to
that device.  The other device will see no output.  If no key is
pressed before the timeout period expires, the system will boot to the
default GRUB boot entry.

Please refer to the GRUB documentation for further information.


\subsubsection{Serial Console Xen configuration}

Enabling Xen serial console output neither enables nor disables Linux
kernel output or logging in to Linux over serial port.  It does however
allow you to monitor and log the Xen boot process via serial console and
can be very useful in debugging.

%% kernel /boot/xen-2.0.gz dom0_mem=131072 com1=115200,8n1
%% module /boot/vmlinuz-2.6-xen0 root=/dev/sda4 ro

In order to configure Xen serial console output, it is necessary to
add a boot option to your GRUB config; e.g.\ replace the previous
example kernel line with:
\begin{quote} {\small \begin{verbatim}
   kernel /boot/xen.gz dom0_mem=131072 com1=115200,8n1
\end{verbatim}}
\end{quote}

This configures Xen to output on COM1 at 115,200 baud, 8 data bits, 1
stop bit and no parity. Modify these parameters for your environment.

One can also configure XenLinux to share the serial console; to achieve
this append ``\path{console=ttyS0}'' to your module line.


\subsubsection{Serial Console Linux configuration}

Enabling Linux serial console output at boot neither enables nor
disables logging in to Linux over serial port.  It does however allow
you to monitor and log the Linux boot process via serial console and can be
very useful in debugging.

To enable Linux output at boot time, add the parameter
\path{console=ttyS0} (or ttyS1, ttyS2, etc.) to your kernel GRUB line.
Under Xen, this might be:
\begin{quote} 
{\footnotesize \begin{verbatim}
  module /vmlinuz-2.6-xen0 ro root=/dev/VolGroup00/LogVol00 \
  console=ttyS0, 115200
\end{verbatim}}
\end{quote}
to enable output over ttyS0 at 115200 baud.



\subsubsection{Serial Console Login configuration}

Logging in to Linux via serial console, under Xen or otherwise, requires
specifying a login prompt be started on the serial port.  To permit root
logins over serial console, the serial port must be added to
\path{/etc/securetty}.

\newpage
To automatically start a login prompt over the serial port, 
add the line: \begin{quote} {\small {\tt c:2345:respawn:/sbin/mingetty
ttyS0}} \end{quote} to \path{/etc/inittab}.   Run \path{init q} to force
a reload of your inttab and start getty.

To enable root logins, add \path{ttyS0} to \path{/etc/securetty} if not
already present.

Your distribution may use an alternate getty; options include getty,
mgetty and agetty.  Consult your distribution's documentation
for further information.


\subsection{TLS Libraries}

Users of the XenLinux 2.6 kernel should disable Thread Local Storage
(TLS) (e.g.\ by doing a \path{mv /lib/tls /lib/tls.disabled}) before
attempting to boot a XenLinux kernel\footnote{If you boot without first
  disabling TLS, you will get a warning message during the boot process.
  In this case, simply perform the rename after the machine is up and
  then run \path{/sbin/ldconfig} to make it take effect.}. You can
always reenable TLS by restoring the directory to its original location
(i.e.\ \path{mv /lib/tls.disabled /lib/tls}).

The reason for this is that the current TLS implementation uses
segmentation in a way that is not permissible under Xen. If TLS is not
disabled, an emulation mode is used within Xen which reduces performance
substantially. To ensure full performance you should install a 
`Xen-friendly' (nosegneg) version of the library. 


\section{Booting Xen}

It should now be possible to restart the system and use Xen. Reboot and
choose the new Xen option when the Grub screen appears.

What follows should look much like a conventional Linux boot. The first
portion of the output comes from Xen itself, supplying low level
information about itself and the underlying hardware. The last portion
of the output comes from XenLinux.

You may see some error messages during the XenLinux boot. These are not
necessarily anything to worry about---they may result from kernel
configuration differences between your XenLinux kernel and the one you
usually use.

When the boot completes, you should be able to log into your system as
usual. If you are unable to log in, you should still be able to reboot
with your normal Linux kernel by selecting it at the GRUB prompt.


% Booting Xen
\chapter{Booting a Xen System}

Booting the system into Xen will bring you up into the privileged
management domain, Domain0. At that point you are ready to create
guest domains and ``boot'' them using the \texttt{xm create} command.

\section{Booting Domain0}

After installation and configuration is complete, reboot the system
and and choose the new Xen option when the Grub screen appears.

What follows should look much like a conventional Linux boot.  The
first portion of the output comes from Xen itself, supplying low level
information about itself and the underlying hardware.  The last
portion of the output comes from XenLinux.

%% KMSelf Wed Nov 30 18:09:37 PST 2005:  We should specify what these are.

When the boot completes, you should be able to log into your system as
usual.  If you are unable to log in, you should still be able to
reboot with your normal Linux kernel by selecting it at the GRUB prompt.

The first step in creating a new domain is to prepare a root
filesystem for it to boot.  Typically, this might be stored in a normal
partition, an LVM or other volume manager partition, a disk file or on
an NFS server.  A simple way to do this is simply to boot from your
standard OS install CD and install the distribution into another
partition on your hard drive.

To start the \xend\ control daemon, type
\begin{quote}
  \verb!# xend start!
\end{quote}

If you wish the daemon to start automatically, see the instructions in
Section~\ref{s:xend}. Once the daemon is running, you can use the
\path{xm} tool to monitor and maintain the domains running on your
system. This chapter provides only a brief tutorial. We provide full
details of the \path{xm} tool in the next chapter.

% \section{From the web interface}
%
% Boot the Xen machine and start Xensv (see Chapter~\ref{cha:xensv}
% for more details) using the command: \\
% \verb_# xensv start_ \\
% This will also start Xend (see Chapter~\ref{cha:xend} for more
% information).
%
% The domain management interface will then be available at {\tt
%   http://your\_machine:8080/}.  This provides a user friendly wizard
% for starting domains and functions for managing running domains.
%
% \section{From the command line}
\section{Booting Guest Domains}

\subsection{Creating a Domain Configuration File}

Before you can start an additional domain, you must create a
configuration file. We provide two example files which you can use as
a starting point:
\begin{itemize}
\item \path{/etc/xen/xmexample1} is a simple template configuration
  file for describing a single VM\@.
\item \path{/etc/xen/xmexample2} file is a template description that
  is intended to be reused for multiple virtual machines.  Setting the
  value of the \path{vmid} variable on the \path{xm} command line
  fills in parts of this template.
\end{itemize}

There are also a number of other examples which you may find useful.
Copy one of these files and edit it as appropriate.  Typical values
you may wish to edit include:

\begin{quote}
\begin{description}
\item[kernel] Set this to the path of the kernel you compiled for use
  with Xen (e.g.\ \path{kernel = ``/boot/vmlinuz-2.6-xenU''})
\item[memory] Set this to the size of the domain's memory in megabytes
  (e.g.\ \path{memory = 64})
\item[disk] Set the first entry in this list to calculate the offset
  of the domain's root partition, based on the domain ID\@.  Set the
  second to the location of \path{/usr} if you are sharing it between
  domains (e.g.\ \path{disk = ['phy:your\_hard\_drive\%d,sda1,w' \%
    (base\_partition\_number + vmid),
    'phy:your\_usr\_partition,sda6,r' ]}
\item[dhcp] Uncomment the dhcp variable, so that the domain will
  receive its IP address from a DHCP server (e.g.\ \path{dhcp=``dhcp''})
\end{description}
\end{quote}

You may also want to edit the {\bf vif} variable in order to choose
the MAC address of the virtual ethernet interface yourself.  For
example:

\begin{quote}
\verb_vif = ['mac=00:16:3E:F6:BB:B3']_
\end{quote}
If you do not set this variable, \xend\ will automatically generate a
random MAC address from the range 00:16:3E:xx:xx:xx, assigned by IEEE to
XenSource as an OUI (organizationally unique identifier).  XenSource
Inc. gives permission for anyone to use addresses randomly allocated
from this range for use by their Xen domains.

For a list of IEEE OUI assignments, see 
\url{http://standards.ieee.org/regauth/oui/oui.txt} 


\subsection{Booting the Guest Domain}

The \path{xm} tool provides a variety of commands for managing
domains.  Use the \path{create} command to start new domains. Assuming
you've created a configuration file \path{myvmconf} based around
\path{/etc/xen/xmexample2}, to start a domain with virtual machine
ID~1 you should type:

\begin{quote}
\begin{verbatim}
# xm create -c myvmconf vmid=1
\end{verbatim}
\end{quote}

The \path{-c} switch causes \path{xm} to turn into the domain's
console after creation.  The \path{vmid=1} sets the \path{vmid}
variable used in the \path{myvmconf} file.

You should see the console boot messages from the new domain appearing
in the terminal in which you typed the command, culminating in a login
prompt.


\section{Starting / Stopping Domains Automatically}

It is possible to have certain domains start automatically at boot
time and to have dom0 wait for all running domains to shutdown before
it shuts down the system.

To specify a domain is to start at boot-time, place its configuration
file (or a link to it) under \path{/etc/xen/auto/}.

A Sys-V style init script for Red Hat and LSB-compliant systems is
provided and will be automatically copied to \path{/etc/init.d/}
during install.  You can then enable it in the appropriate way for
your distribution.

For instance, on Red Hat:

\begin{quote}
  \verb_# chkconfig --add xendomains_
\end{quote}

By default, this will start the boot-time domains in runlevels 3, 4
and 5.

You can also use the \path{service} command to run this script
manually, e.g:

\begin{quote}
  \verb_# service xendomains start_

  Starts all the domains with config files under /etc/xen/auto/.
\end{quote}

\begin{quote}
  \verb_# service xendomains stop_

  Shuts down all running Xen domains.
\end{quote}



\part{Configuration and Management}

%% Chapter Domain Management Tools and Daemons
\chapter{Domain Management Tools}

This chapter summarizes the management software and tools available.


\section{\Xend\ }
\label{s:xend}


The \Xend\ node control daemon performs system management functions
related to virtual machines. It forms a central point of control of
virtualized resources, and must be running in order to start and manage
virtual machines. \Xend\ must be run as root because it needs access to
privileged system management functions.

An initialization script named \texttt{/etc/init.d/xend} is provided to
start \Xend\ at boot time. Use the tool appropriate (i.e. chkconfig) for
your Linux distribution to specify the runlevels at which this script
should be executed, or manually create symbolic links in the correct
runlevel directories.

\Xend\ can be started on the command line as well, and supports the
following set of parameters:

\begin{tabular}{ll}
  \verb!# xend start! & start \xend, if not already running \\
  \verb!# xend stop!  & stop \xend\ if already running       \\
  \verb!# xend restart! & restart \xend\ if running, otherwise start it \\
  % \verb!# xend trace_start! & start \xend, with very detailed debug logging \\
  \verb!# xend status! & indicates \xend\ status by its return code
\end{tabular}

A SysV init script called {\tt xend} is provided to start \xend\ at
boot time. {\tt make install} installs this script in
\path{/etc/init.d}. To enable it, you have to make symbolic links in
the appropriate runlevel directories or use the {\tt chkconfig} tool,
where available.  Once \xend\ is running, administration can be done
using the \texttt{xm} tool.

\subsection{Logging}

As \xend\ runs, events will be logged to \path{/var/log/xend.log} and
(less frequently) to \path{/var/log/xend-debug.log}. These, along with
the standard syslog files, are useful when troubleshooting problems.

\subsection{Configuring \Xend\ }

\Xend\ is written in Python. At startup, it reads its configuration
information from the file \path{/etc/xen/xend-config.sxp}. The Xen
installation places an example \texttt{xend-config.sxp} file in the
\texttt{/etc/xen} subdirectory which should work for most installations.

See the example configuration file \texttt{xend-debug.sxp} and the
section 5 man page \texttt{xend-config.sxp} for a full list of
parameters and more detailed information. Some of the most important
parameters are discussed below.

An HTTP interface and a Unix domain socket API are available to
communicate with \Xend. This allows remote users to pass commands to the
daemon. By default, \Xend does not start an HTTP server. It does start a
Unix domain socket management server, as the low level utility
\texttt{xm} requires it. For support of cross-machine migration, \Xend\
can start a relocation server. This support is not enabled by default
for security reasons.

Note: the example \texttt{xend} configuration file modifies the defaults and
starts up \Xend\ as an HTTP server as well as a relocation server.

From the file:

\begin{verbatim}
#(xend-http-server no)
(xend-http-server yes)
#(xend-unix-server yes)
#(xend-relocation-server no)
(xend-relocation-server yes)
\end{verbatim}

Comment or uncomment lines in that file to disable or enable features
that you require.

Connections from remote hosts are disabled by default:

\begin{verbatim}
# Address xend should listen on for HTTP connections, if xend-http-server is
# set.
# Specifying 'localhost' prevents remote connections.
# Specifying the empty string '' (the default) allows all connections.
#(xend-address '')
(xend-address localhost)
\end{verbatim}

It is recommended that if migration support is not needed, the
\texttt{xend-relocation-server} parameter value be changed to
``\texttt{no}'' or commented out.

\section{Xm}
\label{s:xm}

The xm tool is the primary tool for managing Xen from the console. The
general format of an xm command line is:

\begin{verbatim}
# xm command [switches] [arguments] [variables]
\end{verbatim}

The available \emph{switches} and \emph{arguments} are dependent on the
\emph{command} chosen. The \emph{variables} may be set using
declarations of the form {\tt variable=value} and command line
declarations override any of the values in the configuration file being
used, including the standard variables described above and any custom
variables (for instance, the \path{xmdefconfig} file uses a {\tt vmid}
variable).

For online help for the commands available, type:

\begin{quote}
\begin{verbatim}
# xm help
\end{verbatim}
\end{quote}

This will list the most commonly used commands.  The full list can be obtained
using \verb_xm help --long_.  You can also type \path{xm help $<$command$>$}
for more information on a given command.

\subsection{Basic Management Commands}

One useful command is \verb_# xm list_ which lists all domains running in rows
of the following format:
\begin{center} {\tt name domid memory vcpus state cputime}
\end{center}

The meaning of each field is as follows: 
\begin{quote}
  \begin{description}
  \item[name] The descriptive name of the virtual machine.
  \item[domid] The number of the domain ID this virtual machine is
    running in.
  \item[memory] Memory size in megabytes.
  \item[vcpus] The number of virtual CPUs this domain has.
  \item[state] Domain state consists of 5 fields:
    \begin{description}
    \item[r] running
    \item[b] blocked
    \item[p] paused
    \item[s] shutdown
    \item[c] crashed
    \end{description}
  \item[cputime] How much CPU time (in seconds) the domain has used so
    far.
  \end{description}
\end{quote}

The \path{xm list} command also supports a long output format when the
\path{-l} switch is used.  This outputs the full details of the
running domains in \xend's SXP configuration format.


You can get access to the console of a particular domain using 
the \verb_# xm console_ command  (e.g.\ \verb_# xm console myVM_). 



%% Chapter Domain Configuration
\chapter{Domain Configuration}
\label{cha:config}

The following contains the syntax of the domain configuration files
and description of how to further specify networking, driver domain
and general scheduling behavior.


\section{Configuration Files}
\label{s:cfiles}

Xen configuration files contain the following standard variables.
Unless otherwise stated, configuration items should be enclosed in
quotes: see the configuration scripts in \path{/etc/xen/} 
for concrete examples. 

\begin{description}
\item[kernel] Path to the kernel image.
\item[ramdisk] Path to a ramdisk image (optional).
  % \item[builder] The name of the domain build function (e.g.
  %   {\tt'linux'} or {\tt'netbsd'}.
\item[memory] Memory size in megabytes.
\item[vcpus] The number of virtual CPUs. 
\item[console] Port to export the domain console on (default 9600 +
  domain ID).
\item[vif] Network interface configuration.  This may simply contain
an empty string for each desired interface, or may override various
settings, e.g.\ 
\begin{verbatim}
vif = [ 'mac=00:16:3E:00:00:11, bridge=xen-br0',
        'bridge=xen-br1' ]
\end{verbatim}
  to assign a MAC address and bridge to the first interface and assign
  a different bridge to the second interface, leaving \xend\ to choose
  the MAC address.  The settings that may be overridden in this way are
  type, mac, bridge, ip, script, backend, and vifname.
\item[disk] List of block devices to export to the domain e.g. 
  \verb_disk = [ 'phy:hda1,sda1,r' ]_ 
  exports physical device \path{/dev/hda1} to the domain as
  \path{/dev/sda1} with read-only access. Exporting a disk read-write
  which is currently mounted is dangerous -- if you are \emph{certain}
  you wish to do this, you can specify \path{w!} as the mode.
\item[dhcp] Set to {\tt `dhcp'} if you want to use DHCP to configure
  networking.
\item[netmask] Manually configured IP netmask.
\item[gateway] Manually configured IP gateway.
\item[hostname] Set the hostname for the virtual machine.
\item[root] Specify the root device parameter on the kernel command
  line.
\item[nfs\_server] IP address for the NFS server (if any).
\item[nfs\_root] Path of the root filesystem on the NFS server (if
  any).
\item[extra] Extra string to append to the kernel command line (if
  any)
\end{description}

Additional fields are documented in the example configuration files 
(e.g. to configure virtual TPM functionality). 

For additional flexibility, it is also possible to include Python
scripting commands in configuration files.  An example of this is the
\path{xmexample2} file, which uses Python code to handle the
\path{vmid} variable.


%\part{Advanced Topics}


\section{Network Configuration}

For many users, the default installation should work ``out of the
box''.  More complicated network setups, for instance with multiple
Ethernet interfaces and/or existing bridging setups will require some
special configuration.

The purpose of this section is to describe the mechanisms provided by
\xend\ to allow a flexible configuration for Xen's virtual networking.

\subsection{Xen virtual network topology}

Each domain network interface is connected to a virtual network
interface in dom0 by a point to point link (effectively a ``virtual
crossover cable'').  These devices are named {\tt
  vif$<$domid$>$.$<$vifid$>$} (e.g.\ {\tt vif1.0} for the first
interface in domain~1, {\tt vif3.1} for the second interface in
domain~3).

Traffic on these virtual interfaces is handled in domain~0 using
standard Linux mechanisms for bridging, routing, rate limiting, etc.
Xend calls on two shell scripts to perform initial configuration of
the network and configuration of new virtual interfaces.  By default,
these scripts configure a single bridge for all the virtual
interfaces.  Arbitrary routing / bridging configurations can be
configured by customizing the scripts, as described in the following
section.

\subsection{Xen networking scripts}

Xen's virtual networking is configured by two shell scripts (by
default \path{network-bridge} and \path{vif-bridge}).  These are called
automatically by \xend\ when certain events occur, with arguments to
the scripts providing further contextual information.  These scripts
are found by default in \path{/etc/xen/scripts}.  The names and
locations of the scripts can be configured in
\path{/etc/xen/xend-config.sxp}.

\begin{description}
\item[network-bridge:] This script is called whenever \xend\ is started or
  stopped to respectively initialize or tear down the Xen virtual
  network. In the default configuration initialization creates the
  bridge `xen-br0' and moves eth0 onto that bridge, modifying the
  routing accordingly. When \xend\ exits, it deletes the Xen bridge
  and removes eth0, restoring the normal IP and routing configuration.

  %% In configurations where the bridge already exists, this script
  %% could be replaced with a link to \path{/bin/true} (for instance).

\item[vif-bridge:] This script is called for every domain virtual
  interface and can configure firewalling rules and add the vif to the
  appropriate bridge. By default, this adds and removes VIFs on the
  default Xen bridge.
\end{description}

Other example scripts are available (\path{network-route} and
\path{vif-route}, \path{network-nat} and \path{vif-nat}).
For more complex network setups (e.g.\ where routing is required or
integrate with existing bridges) these scripts may be replaced with
customized variants for your site's preferred configuration.

\section{Driver Domain Configuration}
\label{s:ddconf}

\subsection{PCI}
\label{ss:pcidd}

Individual PCI devices can be assigned to a given domain to allow that
domain direct access to the PCI hardware. To use this functionality, ensure
that the PCI Backend is compiled in to a privileged domain (e.g. domain 0)
and that the domains which will be assigned PCI devices have the PCI Frontend
compiled in. In XenLinux, the PCI Backend is available under the Xen
configuration section while the PCI Frontend is under the
architecture-specific "Bus Options" section. You may compile both the backend
and the frontend into the same kernel; they will not affect each other.

The PCI devices you wish to assign to unprivileged domains must be "hidden"
from your backend domain (usually domain 0) so that it does not load a driver
for them. Use the \path{pciback.hide} kernel parameter which is specified on
the kernel command-line and is configurable through GRUB (see
Section~\ref{s:configure}). Note that devices are not really hidden from the
backend domain. The PCI Backend ensures that no other device driver loads
for those devices. PCI devices are identified by hexadecimal
slot/funciton numbers (on Linux, use \path{lspci} to determine slot/funciton
numbers of your devices) and can be specified with or without the PCI domain: \\
\centerline{  {\tt ({\em bus}:{\em slot}.{\em func})} example {\tt (02:1d.3)}} \\
\centerline{  {\tt ({\em domain}:{\em bus}:{\em slot}.{\em func})} example {\tt (0000:02:1d.3)}} \\

An example kernel command-line which hides two PCI devices might be: \\
\centerline{ {\tt root=/dev/sda4 ro console=tty0 pciback.hide=(02:01.f)(0000:04:1d.0) } } \\

To configure a domU to receive a PCI device:

\begin{description}
\item[Command-line:]
  Use the {\em pci} command-line flag. For multiple devices, use the option
  multiple times. \\
\centerline{  {\tt xm create netcard-dd pci=01:00.0 pci=02:03.0 }} \\

\item[Flat Format configuration file:]
  Specify all of your PCI devices in a python list named {\em pci}. \\
\centerline{  {\tt pci=['01:00.0','02:03.0'] }} \\

\item[SXP Format configuration file:]
  Use a single PCI device section for all of your devices (specify the numbers
  in hexadecimal with the preceding '0x'). Note that {\em domain} here refers
  to the PCI domain, not a virtual machine within Xen.
{\small
\begin{verbatim}
(device (pci
    (dev (domain 0x0)(bus 0x3)(slot 0x1a)(func 0x1)
    (dev (domain 0x0)(bus 0x1)(slot 0x5)(func 0x0)
)
\end{verbatim}
}
\end{description}

There are a number of security concerns associated with PCI Driver Domains
that you can read about in Section~\ref{s:ddsecurity}.

%% There are two possible types of privileges: IO privileges and
%% administration privileges.




% Chapter Storage and FileSytem Management
\chapter{Storage and File System Management}

Storage can be made available to virtual machines in a number of
different ways.  This chapter covers some possible configurations.

The most straightforward method is to export a physical block device (a
hard drive or partition) from dom0 directly to the guest domain as a
virtual block device (VBD).

Storage may also be exported from a filesystem image or a partitioned
filesystem image as a \emph{file-backed VBD}.

Finally, standard network storage protocols such as NBD, iSCSI, NFS,
etc., can be used to provide storage to virtual machines.


\section{Exporting Physical Devices as VBDs}
\label{s:exporting-physical-devices-as-vbds}

One of the simplest configurations is to directly export individual
partitions from domain~0 to other domains. To achieve this use the
\path{phy:} specifier in your domain configuration file. For example a
line like
\begin{quote}
  \verb_disk = ['phy:hda3,sda1,w']_
\end{quote}
specifies that the partition \path{/dev/hda3} in domain~0 should be
exported read-write to the new domain as \path{/dev/sda1}; one could
equally well export it as \path{/dev/hda} or \path{/dev/sdb5} should
one wish.

In addition to local disks and partitions, it is possible to export
any device that Linux considers to be ``a disk'' in the same manner.
For example, if you have iSCSI disks or GNBD volumes imported into
domain~0 you can export these to other domains using the \path{phy:}
disk syntax. E.g.:
\begin{quote}
  \verb_disk = ['phy:vg/lvm1,sda2,w']_
\end{quote}

\begin{center}
  \framebox{\bf Warning: Block device sharing}
\end{center}
\begin{quote}
  Block devices should typically only be shared between domains in a
  read-only fashion otherwise the Linux kernel's file systems will get
  very confused as the file system structure may change underneath
  them (having the same ext3 partition mounted \path{rw} twice is a
  sure fire way to cause irreparable damage)!  \Xend\ will attempt to
  prevent you from doing this by checking that the device is not
  mounted read-write in domain~0, and hasn't already been exported
  read-write to another domain.  If you want read-write sharing,
  export the directory to other domains via NFS from domain~0 (or use
  a cluster file system such as GFS or ocfs2).
\end{quote}


\section{Using File-backed VBDs}

It is also possible to use a file in Domain~0 as the primary storage
for a virtual machine.  As well as being convenient, this also has the
advantage that the virtual block device will be \emph{sparse} ---
space will only really be allocated as parts of the file are used.  So
if a virtual machine uses only half of its disk space then the file
really takes up half of the size allocated.

For example, to create a 2GB sparse file-backed virtual block device
(actually only consumes 1KB of disk):
\begin{quote}
  \verb_# dd if=/dev/zero of=vm1disk bs=1k seek=2048k count=1_
\end{quote}

Make a file system in the disk file:
\begin{quote}
  \verb_# mkfs -t ext3 vm1disk_
\end{quote}

(when the tool asks for confirmation, answer `y')

Populate the file system e.g.\ by copying from the current root:
\begin{quote}
\begin{verbatim}
# mount -o loop vm1disk /mnt
# cp -ax /{root,dev,var,etc,usr,bin,sbin,lib} /mnt
# mkdir /mnt/{proc,sys,home,tmp}
\end{verbatim}
\end{quote}

Tailor the file system by editing \path{/etc/fstab},
\path{/etc/hostname}, etc.\ Don't forget to edit the files in the
mounted file system, instead of your domain~0 filesystem, e.g.\ you
would edit \path{/mnt/etc/fstab} instead of \path{/etc/fstab}.  For
this example put \path{/dev/sda1} to root in fstab.

Now unmount (this is important!):
\begin{quote}
  \verb_# umount /mnt_
\end{quote}

In the configuration file set:
\begin{quote}
  \verb_disk = ['file:/full/path/to/vm1disk,sda1,w']_
\end{quote}

As the virtual machine writes to its `disk', the sparse file will be
filled in and consume more space up to the original 2GB.

{\bf Note that file-backed VBDs may not be appropriate for backing
  I/O-intensive domains.}  File-backed VBDs are known to experience
substantial slowdowns under heavy I/O workloads, due to the I/O
handling by the loopback block device used to support file-backed VBDs
in dom0.  Better I/O performance can be achieved by using either
LVM-backed VBDs (Section~\ref{s:using-lvm-backed-vbds}) or physical
devices as VBDs (Section~\ref{s:exporting-physical-devices-as-vbds}).

Linux supports a maximum of eight file-backed VBDs across all domains
by default.  This limit can be statically increased by using the
\emph{max\_loop} module parameter if CONFIG\_BLK\_DEV\_LOOP is
compiled as a module in the dom0 kernel, or by using the
\emph{max\_loop=n} boot option if CONFIG\_BLK\_DEV\_LOOP is compiled
directly into the dom0 kernel.


\section{Using LVM-backed VBDs}
\label{s:using-lvm-backed-vbds}

A particularly appealing solution is to use LVM volumes as backing for
domain file-systems since this allows dynamic growing/shrinking of
volumes as well as snapshot and other features.

To initialize a partition to support LVM volumes:
\begin{quote}
\begin{verbatim}
# pvcreate /dev/sda10           
\end{verbatim} 
\end{quote}

Create a volume group named `vg' on the physical partition:
\begin{quote}
\begin{verbatim}
# vgcreate vg /dev/sda10
\end{verbatim} 
\end{quote}

Create a logical volume of size 4GB named `myvmdisk1':
\begin{quote}
\begin{verbatim}
# lvcreate -L4096M -n myvmdisk1 vg
\end{verbatim}
\end{quote}

You should now see that you have a \path{/dev/vg/myvmdisk1} Make a
filesystem, mount it and populate it, e.g.:
\begin{quote}
\begin{verbatim}
# mkfs -t ext3 /dev/vg/myvmdisk1
# mount /dev/vg/myvmdisk1 /mnt
# cp -ax / /mnt
# umount /mnt
\end{verbatim}
\end{quote}

Now configure your VM with the following disk configuration:
\begin{quote}
\begin{verbatim}
 disk = [ 'phy:vg/myvmdisk1,sda1,w' ]
\end{verbatim}
\end{quote}

LVM enables you to grow the size of logical volumes, but you'll need
to resize the corresponding file system to make use of the new space.
Some file systems (e.g.\ ext3) now support online resize.  See the LVM
manuals for more details.

You can also use LVM for creating copy-on-write (CoW) clones of LVM
volumes (known as writable persistent snapshots in LVM terminology).
This facility is new in Linux 2.6.8, so isn't as stable as one might
hope.  In particular, using lots of CoW LVM disks consumes a lot of
dom0 memory, and error conditions such as running out of disk space
are not handled well. Hopefully this will improve in future.

To create two copy-on-write clones of the above file system you would
use the following commands:

\begin{quote}
\begin{verbatim}
# lvcreate -s -L1024M -n myclonedisk1 /dev/vg/myvmdisk1
# lvcreate -s -L1024M -n myclonedisk2 /dev/vg/myvmdisk1
\end{verbatim}
\end{quote}

Each of these can grow to have 1GB of differences from the master
volume. You can grow the amount of space for storing the differences
using the lvextend command, e.g.:
\begin{quote}
\begin{verbatim}
# lvextend +100M /dev/vg/myclonedisk1
\end{verbatim}
\end{quote}

Don't let the `differences volume' ever fill up otherwise LVM gets
rather confused. It may be possible to automate the growing process by
using \path{dmsetup wait} to spot the volume getting full and then
issue an \path{lvextend}.

In principle, it is possible to continue writing to the volume that
has been cloned (the changes will not be visible to the clones), but
we wouldn't recommend this: have the cloned volume as a `pristine'
file system install that isn't mounted directly by any of the virtual
machines.


\section{Using NFS Root}

First, populate a root filesystem in a directory on the server
machine. This can be on a distinct physical machine, or simply run
within a virtual machine on the same node.

Now configure the NFS server to export this filesystem over the
network by adding a line to \path{/etc/exports}, for instance:

\begin{quote}
  \begin{small}
\begin{verbatim}
/export/vm1root      1.2.3.4/24 (rw,sync,no_root_squash)
\end{verbatim}
  \end{small}
\end{quote}

Finally, configure the domain to use NFS root.  In addition to the
normal variables, you should make sure to set the following values in
the domain's configuration file:

\begin{quote}
  \begin{small}
\begin{verbatim}
root       = '/dev/nfs'
nfs_server = '2.3.4.5'       # substitute IP address of server
nfs_root   = '/path/to/root' # path to root FS on the server
\end{verbatim}
  \end{small}
\end{quote}

The domain will need network access at boot time, so either statically
configure an IP address using the config variables \path{ip},
\path{netmask}, \path{gateway}, \path{hostname}; or enable DHCP
(\path{dhcp='dhcp'}).

Note that the Linux NFS root implementation is known to have stability
problems under high load (this is not a Xen-specific problem), so this
configuration may not be appropriate for critical servers.


\chapter{CPU Management}

%% KMS Something sage about CPU / processor management.

Xen allows a domain's virtual CPU(s) to be associated with one or more
host CPUs.  This can be used to allocate real resources among one or
more guests, or to make optimal use of processor resources when
utilizing dual-core, hyperthreading, or other advanced CPU technologies.

Xen enumerates physical CPUs in a `depth first' fashion.  For a system
with both hyperthreading and multiple cores, this would be all the
hyperthreads on a given core, then all the cores on a given socket,
and then all sockets.  I.e.  if you had a two socket, dual core,
hyperthreaded Xeon the CPU order would be:


\begin{center}
\begin{tabular}{l|l|l|l|l|l|l|r}
\multicolumn{4}{c|}{socket0}     &  \multicolumn{4}{c}{socket1} \\ \hline
\multicolumn{2}{c|}{core0}  &  \multicolumn{2}{c|}{core1}  &
\multicolumn{2}{c|}{core0}  &  \multicolumn{2}{c}{core1} \\ \hline
ht0 & ht1 & ht0 & ht1 & ht0 & ht1 & ht0 & ht1 \\
\#0 & \#1 & \#2 & \#3 & \#4 & \#5 & \#6 & \#7 \\
\end{tabular}
\end{center}


Having multiple vcpus belonging to the same domain mapped to the same
physical CPU is very likely to lead to poor performance. It's better to
use `vcpus-set' to hot-unplug one of the vcpus and ensure the others are
pinned on different CPUs.

If you are running IO intensive tasks, its typically better to dedicate
either a hyperthread or whole core to running domain 0, and hence pin
other domains so that they can't use CPU 0. If your workload is mostly
compute intensive, you may want to pin vcpus such that all physical CPU
threads are available for guest domains.

\chapter{Migrating Domains}

\section{Domain Save and Restore}

The administrator of a Xen system may suspend a virtual machine's
current state into a disk file in domain~0, allowing it to be resumed at
a later time.

For example you can suspend a domain called ``VM1'' to disk using the
command:
\begin{verbatim}
# xm save VM1 VM1.chk
\end{verbatim}

This will stop the domain named ``VM1'' and save its current state
into a file called \path{VM1.chk}.

To resume execution of this domain, use the \path{xm restore} command:
\begin{verbatim}
# xm restore VM1.chk
\end{verbatim}

This will restore the state of the domain and resume its execution.
The domain will carry on as before and the console may be reconnected
using the \path{xm console} command, as described earlier.

\section{Migration and Live Migration}

Migration is used to transfer a domain between physical hosts. There
are two varieties: regular and live migration. The former moves a
virtual machine from one host to another by pausing it, copying its
memory contents, and then resuming it on the destination. The latter
performs the same logical functionality but without needing to pause
the domain for the duration. In general when performing live migration
the domain continues its usual activities and---from the user's
perspective---the migration should be imperceptible.

To perform a live migration, both hosts must be running Xen / \xend\ and
the destination host must have sufficient resources (e.g.\ memory
capacity) to accommodate the domain after the move. Furthermore we
currently require both source and destination machines to be on the same
L2 subnet.

Currently, there is no support for providing automatic remote access
to filesystems stored on local disk when a domain is migrated.
Administrators should choose an appropriate storage solution (i.e.\
SAN, NAS, etc.) to ensure that domain filesystems are also available
on their destination node. GNBD is a good method for exporting a
volume from one machine to another. iSCSI can do a similar job, but is
more complex to set up.

When a domain migrates, it's MAC and IP address move with it, thus it is
only possible to migrate VMs within the same layer-2 network and IP
subnet. If the destination node is on a different subnet, the
administrator would need to manually configure a suitable etherip or IP
tunnel in the domain~0 of the remote node.

A domain may be migrated using the \path{xm migrate} command. To live
migrate a domain to another machine, we would use the command:

\begin{verbatim}
# xm migrate --live mydomain destination.ournetwork.com
\end{verbatim}

Without the \path{--live} flag, \xend\ simply stops the domain and
copies the memory image over to the new node and restarts it. Since
domains can have large allocations this can be quite time consuming,
even on a Gigabit network. With the \path{--live} flag \xend\ attempts
to keep the domain running while the migration is in progress, resulting
in typical down times of just 60--300ms.

For now it will be necessary to reconnect to the domain's console on the
new machine using the \path{xm console} command. If a migrated domain
has any open network connections then they will be preserved, so SSH
connections do not have this limitation.


%% Chapter Securing Xen
\chapter{Securing Xen}

This chapter describes how to secure a Xen system. It describes a number
of scenarios and provides a corresponding set of best practices. It
begins with a section devoted to understanding the security implications
of a Xen system.


\section{Xen Security Considerations}

When deploying a Xen system, one must be sure to secure the management
domain (Domain-0) as much as possible. If the management domain is
compromised, all other domains are also vulnerable. The following are a
set of best practices for Domain-0:

\begin{enumerate}
\item \textbf{Run the smallest number of necessary services.} The less
  things that are present in a management partition, the better.
  Remember, a service running as root in the management domain has full
  access to all other domains on the system.
\item \textbf{Use a firewall to restrict the traffic to the management
    domain.} A firewall with default-reject rules will help prevent
  attacks on the management domain.
\item \textbf{Do not allow users to access Domain-0.} The Linux kernel
  has been known to have local-user root exploits. If you allow normal
  users to access Domain-0 (even as unprivileged users) you run the risk
  of a kernel exploit making all of your domains vulnerable.
\end{enumerate}

\section{Driver Domain Security Considerations}
\label{s:ddsecurity}

Driver domains address a range of security problems that exist regarding
the use of device drivers and hardware. On many operating systems in common
use today, device drivers run within the kernel with the same privileges as
the kernel. Few or no mechanisms exist to protect the integrity of the kernel
from a misbehaving (read "buggy") or malicious device driver. Driver
domains exist to aid in isolating a device driver within its own virtual
machine where it cannot affect the stability and integrity of other
domains. If a driver crashes, the driver domain can be restarted rather than
have the entire machine crash (and restart) with it. Drivers written by
unknown or untrusted third-parties can be confined to an isolated space.
Driver domains thus address a number of security and stability issues with
device drivers.

However, due to limitations in current hardware, a number of security
concerns remain that need to be considered when setting up driver domains (it
should be noted that the following list is not intended to be exhaustive).

\begin{enumerate}
\item \textbf{Without an IOMMU, a hardware device can DMA to memory regions
  outside of its controlling domain.} Architectures which do not have an
  IOMMU (e.g. most x86-based platforms) to restrict DMA usage by hardware
  are vulnerable. A hardware device which can perform arbitrary memory reads
  and writes can read/write outside of the memory of its controlling domain.
  A malicious or misbehaving domain could use a hardware device it controls
  to send data overwriting memory in another domain or to read arbitrary
  regions of memory in another domain.
\item \textbf{Shared buses are vulnerable to sniffing.} Devices that share
  a data bus can sniff (and possible spoof) each others' data. Device A that
  is assigned to Domain A could eavesdrop on data being transmitted by
  Domain B to Device B and then relay that data back to Domain A.
\item \textbf{Devices which share interrupt lines can either prevent the
  reception of that interrupt by the driver domain or can trigger the
  interrupt service routine of that guest needlessly.} A devices which shares
  a level-triggered interrupt (e.g. PCI devices) with another device can
  raise an interrupt and never clear it. This effectively blocks other devices
  which share that interrupt line from notifying their controlling driver
  domains that they need to be serviced. A device which shares an
  any type of interrupt line can trigger its interrupt continually which
  forces execution time to be spent (in multiple guests) in the interrupt
  service routine (potentially denying time to other processes within that
  guest). System architectures which allow each device to have its own
  interrupt line (e.g. PCI's Message Signaled Interrupts) are less
  vulnerable to this denial-of-service problem.
\item \textbf{Devices may share the use of I/O memory address space.} Xen can
  only restrict access to a device's physical I/O resources at a certain
  granularity. For interrupt lines and I/O port address space, that
  granularity is very fine (per interrupt line and per I/O port). However,
  Xen can only restrict access to I/O memory address space on a page size
  basis. If more than one device shares use of a page in I/O memory address
  space, the domains to which those devices are assigned will be able to
  access the I/O memory address space of each other's devices.
\end{enumerate}


\section{Security Scenarios}


\subsection{The Isolated Management Network}

In this scenario, each node has two network cards in the cluster. One
network card is connected to the outside world and one network card is a
physically isolated management network specifically for Xen instances to
use.

As long as all of the management partitions are trusted equally, this is
the most secure scenario. No additional configuration is needed other
than forcing Xend to bind to the management interface for relocation.


\subsection{A Subnet Behind a Firewall}

In this scenario, each node has only one network card but the entire
cluster sits behind a firewall. This firewall should do at least the
following:

\begin{enumerate}
\item Prevent IP spoofing from outside of the subnet.
\item Prevent access to the relocation port of any of the nodes in the
  cluster except from within the cluster.
\end{enumerate}

The following iptables rules can be used on each node to prevent
migrations to that node from outside the subnet assuming the main
firewall does not do this for you:

\begin{verbatim}
# this command disables all access to the Xen relocation
# port:
iptables -A INPUT -p tcp --destination-port 8002 -j REJECT

# this command enables Xen relocations only from the specific
# subnet:
iptables -I INPUT -p tcp -{}-source 192.168.1.1/8 \
    --destination-port 8002 -j ACCEPT
\end{verbatim}

\subsection{Nodes on an Untrusted Subnet}

Migration on an untrusted subnet is not safe in current versions of Xen.
It may be possible to perform migrations through a secure tunnel via an
VPN or SSH. The only safe option in the absence of a secure tunnel is to
disable migration completely. The easiest way to do this is with
iptables:

\begin{verbatim}
# this command disables all access to the Xen relocation port
iptables -A INPUT -p tcp -{}-destination-port 8002 -j REJECT
\end{verbatim}

\part{Reference}

%% Chapter Build and Boot Options
\chapter{Build and Boot Options} 

This chapter describes the build- and boot-time options which may be
used to tailor your Xen system.

\section{Top-level Configuration Options} 

Top-level configuration is achieved by editing one of two 
files: \path{Config.mk} and \path{Makefile}. 

The former allows the overall build target architecture to be 
specified. You will typically not need to modify this unless 
you are cross-compiling or if you wish to build a PAE-enabled 
Xen system. Additional configuration options are documented 
in the \path{Config.mk} file. 

The top-level \path{Makefile} is chiefly used to customize the set of
kernels built. Look for the line: 
\begin{quote}
\begin{verbatim}
KERNELS ?= linux-2.6-xen0 linux-2.6-xenU
\end{verbatim}
\end{quote}

Allowable options here are any kernels which have a corresponding 
build configuration file in the \path{buildconfigs/} directory. 



\section{Xen Build Options}

Xen provides a number of build-time options which should be set as
environment variables or passed on make's command-line.

\begin{description}
\item[verbose=y] Enable debugging messages when Xen detects an
  unexpected condition.  Also enables console output from all domains.
\item[debug=y] Enable debug assertions.  Implies {\bf verbose=y}.
  (Primarily useful for tracing bugs in Xen).
\item[debugger=y] Enable the in-Xen debugger. This can be used to
  debug Xen, guest OSes, and applications.
\item[perfc=y] Enable performance counters for significant events
  within Xen. The counts can be reset or displayed on Xen's console
  via console control keys.
\end{description}


\section{Xen Boot Options}
\label{s:xboot}

These options are used to configure Xen's behaviour at runtime.  They
should be appended to Xen's command line, either manually or by
editing \path{grub.conf}.

\begin{description}
\item [ noreboot ] Don't reboot the machine automatically on errors.
  This is useful to catch debug output if you aren't catching console
  messages via the serial line.
\item [ nosmp ] Disable SMP support.  This option is implied by
  `ignorebiostables'.
\item [ watchdog ] Enable NMI watchdog which can report certain
  failures.
\item [ noirqbalance ] Disable software IRQ balancing and affinity.
  This can be used on systems such as Dell 1850/2850 that have
  workarounds in hardware for IRQ-routing issues.
\item [ badpage=$<$page number$>$,$<$page number$>$, \ldots ] Specify
  a list of pages not to be allocated for use because they contain bad
  bytes. For example, if your memory tester says that byte 0x12345678
  is bad, you would place `badpage=0x12345' on Xen's command line.
\item [ com1=$<$baud$>$,DPS,$<$io\_base$>$,$<$irq$>$
  com2=$<$baud$>$,DPS,$<$io\_base$>$,$<$irq$>$ ] \mbox{}\\
  Xen supports up to two 16550-compatible serial ports.  For example:
  `com1=9600, 8n1, 0x408, 5' maps COM1 to a 9600-baud port, 8 data
  bits, no parity, 1 stop bit, I/O port base 0x408, IRQ 5.  If some
  configuration options are standard (e.g., I/O base and IRQ), then
  only a prefix of the full configuration string need be specified. If
  the baud rate is pre-configured (e.g., by the bootloader) then you
  can specify `auto' in place of a numeric baud rate.
\item [ console=$<$specifier list$>$ ] Specify the destination for Xen
  console I/O.  This is a comma-separated list of, for example:
  \begin{description}
  \item[ vga ] Use VGA console and allow keyboard input.
  \item[ com1 ] Use serial port com1.
  \item[ com2H ] Use serial port com2. Transmitted chars will have the
    MSB set. Received chars must have MSB set.
  \item[ com2L] Use serial port com2. Transmitted chars will have the
    MSB cleared. Received chars must have MSB cleared.
  \end{description}
  The latter two examples allow a single port to be shared by two
  subsystems (e.g.\ console and debugger). Sharing is controlled by
  MSB of each transmitted/received character.  [NB. Default for this
  option is `com1,vga']
\item [ sync\_console ] Force synchronous console output. This is
  useful if you system fails unexpectedly before it has sent all
  available output to the console. In most cases Xen will
  automatically enter synchronous mode when an exceptional event
  occurs, but this option provides a manual fallback.
\item [ conswitch=$<$switch-char$><$auto-switch-char$>$ ] Specify how
  to switch serial-console input between Xen and DOM0. The required
  sequence is CTRL-$<$switch-char$>$ pressed three times. Specifying
  the backtick character disables switching.  The
  $<$auto-switch-char$>$ specifies whether Xen should auto-switch
  input to DOM0 when it boots --- if it is `x' then auto-switching is
  disabled.  Any other value, or omitting the character, enables
  auto-switching.  [NB. Default switch-char is `a'.]
\item [ nmi=xxx ]
  Specify what to do with an NMI parity or I/O error. \\
  `nmi=fatal':  Xen prints a diagnostic and then hangs. \\
  `nmi=dom0':   Inform DOM0 of the NMI. \\
  `nmi=ignore': Ignore the NMI.
\item [ mem=xxx ] Set the physical RAM address limit. Any RAM
  appearing beyond this physical address in the memory map will be
  ignored. This parameter may be specified with a B, K, M or G suffix,
  representing bytes, kilobytes, megabytes and gigabytes respectively.
  The default unit, if no suffix is specified, is kilobytes.
\item [ dom0\_mem=xxx ] Set the amount of memory to be allocated to
  domain0. In Xen 3.x the parameter may be specified with a B, K, M or
  G suffix, representing bytes, kilobytes, megabytes and gigabytes
  respectively; if no suffix is specified, the parameter defaults to
  kilobytes. In previous versions of Xen, suffixes were not supported
  and the value is always interpreted as kilobytes.
\item [ tbuf\_size=xxx ] Set the size of the per-cpu trace buffers, in
  pages (default 1).  Note that the trace buffers are only enabled in
  debug builds.  Most users can ignore this feature completely.
\item [ sched=xxx ] Select the CPU scheduler Xen should use.  The
  current possibilities are `sedf' (default) and `bvt'.
\item [ apic\_verbosity=debug,verbose ] Print more detailed
  information about local APIC and IOAPIC configuration.
\item [ lapic ] Force use of local APIC even when left disabled by
  uniprocessor BIOS.
\item [ nolapic ] Ignore local APIC in a uniprocessor system, even if
  enabled by the BIOS.
\item [ apic=bigsmp,default,es7000,summit ] Specify NUMA platform.
  This can usually be probed automatically.
\end{description}

In addition, the following options may be specified on the Xen command
line. Since domain 0 shares responsibility for booting the platform,
Xen will automatically propagate these options to its command line.
These options are taken from Linux's command-line syntax with
unchanged semantics.

\begin{description}
\item [ acpi=off,force,strict,ht,noirq,\ldots ] Modify how Xen (and
  domain 0) parses the BIOS ACPI tables.
\item [ acpi\_skip\_timer\_override ] Instruct Xen (and domain~0) to
  ignore timer-interrupt override instructions specified by the BIOS
  ACPI tables.
\item [ noapic ] Instruct Xen (and domain~0) to ignore any IOAPICs
  that are present in the system, and instead continue to use the
  legacy PIC.
\end{description} 


\section{XenLinux Boot Options}

In addition to the standard Linux kernel boot options, we support:
\begin{description}
\item[ xencons=xxx ] Specify the device node to which the Xen virtual
  console driver is attached. The following options are supported:
  \begin{center}
    \begin{tabular}{l}
      `xencons=off': disable virtual console \\
      `xencons=tty': attach console to /dev/tty1 (tty0 at boot-time) \\
      `xencons=ttyS': attach console to /dev/ttyS0
    \end{tabular}
\end{center}
The default is ttyS for dom0 and tty for all other domains.
\end{description}


%% Chapter Further Support
\chapter{Further Support}

If you have questions that are not answered by this manual, the
sources of information listed below may be of interest to you.  Note
that bug reports, suggestions and contributions related to the
software (or the documentation) should be sent to the Xen developers'
mailing list (address below).


\section{Other Documentation}

For developers interested in porting operating systems to Xen, the
\emph{Xen Interface Manual} is distributed in the \path{docs/}
directory of the Xen source distribution.


\section{Online References}

The official Xen web site can be found at:
\begin{quote} {\tt http://www.xensource.com}
\end{quote}

This contains links to the latest versions of all online
documentation, including the latest version of the FAQ.

Information regarding Xen is also available at the Xen Wiki at
\begin{quote} {\tt http://wiki.xensource.com/xenwiki/}\end{quote}
The Xen project uses Bugzilla as its bug tracking system. You'll find
the Xen Bugzilla at http://bugzilla.xensource.com/bugzilla/.


\section{Mailing Lists}

There are several mailing lists that are used to discuss Xen related
topics. The most widely relevant are listed below. An official page of
mailing lists and subscription information can be found at \begin{quote}
  {\tt http://lists.xensource.com/} \end{quote}

\begin{description}
\item[xen-devel@lists.xensource.com] Used for development
  discussions and bug reports.  Subscribe at: \\
  {\small {\tt http://lists.xensource.com/xen-devel}}
\item[xen-users@lists.xensource.com] Used for installation and usage
  discussions and requests for help.  Subscribe at: \\
  {\small {\tt http://lists.xensource.com/xen-users}}
\item[xen-announce@lists.xensource.com] Used for announcements only.
  Subscribe at: \\
  {\small {\tt http://lists.xensource.com/xen-announce}}
\item[xen-changelog@lists.xensource.com] Changelog feed
  from the unstable and 2.0 trees - developer oriented.  Subscribe at: \\
  {\small {\tt http://lists.xensource.com/xen-changelog}}
\end{description}



%%%%%%%%%%%%%%%%%%%%%%%%%%%%%%%%%%%%%%%%%%%%%%%%%%%%%%%%%%%%%%%%

\appendix

\chapter{Unmodified (VMX) guest domains in Xen with Intel\textregistered Virtualization Technology (VT)}

Xen supports guest domains running unmodified Guest operating systems using Virtualization Technology (VT) available on recent Intel Processors. More information about the Intel Virtualization Technology implementing Virtual Machine Extensions (VMX) in the processor is available on the Intel website at \\
 {\small {\tt http://www.intel.com/technology/computing/vptech}}

\section{Building Xen with VT support}

The following packages need to be installed in order to build Xen with VT support. Some Linux distributions do not provide these packages by default.

\begin{tabular}{lp{11.0cm}}
{\bfseries Package} & {\bfseries Description} \\

dev86 & The dev86 package provides an assembler and linker for real mode 80x86 instructions. You need to have this package installed in order to build the BIOS code which runs in (virtual) real mode. 

If the dev86 package is not available on the x86\_64 distribution, you can install the i386 version of it. The dev86 rpm package for various distributions can be found at {\scriptsize {\tt http://www.rpmfind.net/linux/rpm2html/search.php?query=dev86\&submit=Search}} \\

LibVNCServer & The unmodified guest's VGA display, keyboard, and mouse are virtualized using the vncserver library provided by this package. You can get the sources of libvncserver from {\small {\tt http://sourceforge.net/projects/libvncserver}}. Build and install the sources on the build system to get the libvncserver library. The 0.8pre version of libvncserver is currently working well with Xen.\\

SDL-devel, SDL & Simple DirectMedia Layer (SDL) is another way of virtualizing the unmodified guest console. It provides an X window for the guest console. 

If the SDL and SDL-devel packages are not installed by default on the build system, they can be obtained from  {\scriptsize {\tt http://www.rpmfind.net/linux/rpm2html/search.php?query=SDL\&amp;submit=Search}}
, {\scriptsize {\tt http://www.rpmfind.net/linux/rpm2html/search.php?query=SDL-devel\&submit=Search}} \\

\end{tabular}

\section{Configuration file for unmodified VMX guests}

The Xen installation includes a sample configuration file, {\small {\tt /etc/xen/xmexample.vmx}}. There are comments describing all the options. In addition to the common options that are the same as those for paravirtualized guest configurations, VMX guest configurations have the following settings:

\begin{tabular}{lp{11.0cm}}

{\bfseries Parameter} & {\bfseries Description} \\

kernel &        The VMX firmware loader, {\small {\tt /usr/lib/xen/boot/vmxloader}}\\

builder &       The domain build function. The VMX domain uses the vmx builder.\\

acpi & Enable VMX guest ACPI, default=0 (disabled)\\

apic & Enable VMX guest APIC, default=0 (disabled)\\

vif     & Optionally defines MAC address and/or bridge for the network interfaces. Random MACs are assigned if not given. {\small {\tt type=ioemu}} means ioemu is used to virtualize the VMX NIC. If no type is specified, vbd is used, as with paravirtualized guests.\\

disk & Defines the disk devices you want the domain to have access to, and what you want them accessible as. If using a physical device as the VMX guest's disk, each disk entry is of the form 

{\small {\tt phy:UNAME,ioemu:DEV,MODE,}}

where UNAME is the device, DEV is the device name the domain will see, and MODE is r for read-only, w for read-write. ioemu means the disk will use ioemu to virtualize the VMX disk. If not adding ioemu, it uses vbd like paravirtualized guests.

If using disk image file, its form should be like 

{\small {\tt file:FILEPATH,ioemu:DEV,MODE}}

If using more than one disk, there should be a comma between each disk entry. For example:

{\scriptsize {\tt disk = ['file:/var/images/image1.img,ioemu:hda,w', 'file:/var/images/image2.img,ioemu:hdb,w']}}\\

cdrom   & Disk image for CD-ROM. The default is {\small {\tt /dev/cdrom}} for Domain0. Inside the VMX domain, the CD-ROM will available as device {\small {\tt /dev/hdc}}. The entry can also point to an ISO file.\\

boot & Boot from floppy (a), hard disk (c) or CD-ROM (d). For example, to boot from CD-ROM, the entry should be:

boot='d'\\

device\_model & The device emulation tool for VMX guests. This parameter should not be changed.\\

sdl &   Enable SDL library for graphics, default = 0 (disabled)\\

vnc &   Enable VNC library for graphics, default = 1 (enabled)\\

vncviewer &     Enable spawning of the vncviewer (only valid when vnc=1), default = 1 (enabled)

If vnc=1 and vncviewer=0, user can use vncviewer to manually connect VMX from remote. For example:

{\small {\tt vncviewer domain0\_IP\_address:VMX\_domain\_id}} \\

ne2000 &        Enable ne2000, default = 0 (disabled; use pcnet)\\

serial &        Enable redirection of VMX serial output to pty device\\

localtime &     Set the real time clock to local time [default=0, that is, set to UTC].\\

enable-audio &  Enable audio support. This is under development.\\

full-screen     & Start in full screen. This is under development.\\

nographic &     Another way to redirect serial output. If enabled, no 'sdl' or 'vnc' can work. Not recommended.\\

\end{tabular}


\section{Creating virtual disks from scratch}
\subsection{Using physical disks}
If you are using a physical disk or physical disk partition, you need to install a Linux OS on the disk first. Then the boot loader should be installed in the correct place. For example {\small {\tt dev/sda}} for booting from the whole disk, or {\small {\tt /dev/sda1}} for booting from partition 1.

\subsection{Using disk image files}
You need to create a large empty disk image file first; then, you need to install a Linux OS onto it. There are two methods you can choose. One is directly installing it using a VMX guest while booting from the OS installation CD-ROM. The other is copying an installed OS into it. The boot loader will also need to be installed.

\subsubsection*{To create the image file:}
The image size should be big enough to accommodate the entire OS. This example assumes the size is 1G (which is probably too small for most OSes).

{\small {\tt \# dd if=/dev/zero of=hd.img bs=1M count=1 seek=1023}}

\subsubsection*{To directly install Linux OS into an image file using a VMX guest:}

Install Xen and create VMX with the original image file with booting from CD-ROM. Then it is just like a normal Linux OS installation. The VMX configuration file should have these two entries before creating:

{\small {\tt cdrom='/dev/cdrom'
boot='d'}}

If this method does not succeed, you can choose the following method of copying an installed Linux OS into an image file.

\subsubsection*{To copy a installed OS into an image file:}
Directly installing is an easier way to make partitions and install an OS in a disk image file. But if you want to create a specific OS in your disk image, then you will most likely want to use this method.

\begin{enumerate}
\item {\bfseries Install a normal Linux OS on the host machine}\\
You can choose any way to install Linux, such as using yum to install Red Hat Linux or YAST to install Novell SuSE Linux. The rest of this example assumes the Linux OS is installed in {\small {\tt /var/guestos/}}.

\item {\bfseries Make the partition table}\\
The image file will be treated as hard disk, so you should make the partition table in the image file. For example:

{\scriptsize {\tt \# losetup /dev/loop0 hd.img\\
\# fdisk -b 512 -C 4096 -H 16 -S 32 /dev/loop0\\
press 'n' to add new partition\\
press 'p' to choose primary partition\\
press '1' to set partition number\\
press "Enter" keys to choose default value of "First Cylinder" parameter.\\
press "Enter" keys to choose default value of "Last Cylinder" parameter.\\
press 'w' to write partition table and exit\\
\# losetup -d /dev/loop0}}

\item {\bfseries Make the file system and install grub}\\
{\scriptsize {\tt \# ln -s /dev/loop0 /dev/loop\\
\# losetup /dev/loop0 hd.img\\
\# losetup -o 16384 /dev/loop1 hd.img\\
\# mkfs.ext3 /dev/loop1\\
\# mount /dev/loop1 /mnt\\
\# mkdir -p /mnt/boot/grub\\
\# cp /boot/grub/stage* /boot/grub/e2fs\_stage1\_5 /mnt/boot/grub\\
\# umount /mnt\\
\# grub\\
grub> device (hd0) /dev/loop\\
grub> root (hd0,0)\\
grub> setup (hd0)\\
grub> quit\\
\# rm /dev/loop\\
\# losetup -d /dev/loop0\\
\# losetup -d /dev/loop1}}

The {\small {\tt losetup}} option {\small {\tt -o 16384}} skips the partition table in the image file. It is the number of sectors times 512. We need {\small {\tt /dev/loop}} because grub is expecting a disk device \emph{name}, where \emph{name} represents the entire disk and \emph{name1} represents the first partition.

\item {\bfseries Copy the OS files to the image}\\ 
If you have Xen installed, you can easily use {\small {\tt lomount}} instead of {\small {\tt losetup}} and {\small {\tt mount}} when coping files to some partitions. {\small {\tt lomount}} just needs the partition information.

{\scriptsize {\tt \# lomount -t ext3 -diskimage hd.img -partition 1 /mnt/guest\\
\# cp -ax /var/guestos/\{root,dev,var,etc,usr,bin,sbin,lib\} /mnt/guest\\
\# mkdir /mnt/guest/\{proc,sys,home,tmp\}}}

\item {\bfseries Edit the {\small {\tt /etc/fstab}} of the guest image}\\
The fstab should look like this:

{\scriptsize {\tt \# vim /mnt/guest/etc/fstab\\
/dev/hda1       /               ext3            defaults 1 1\\
none            /dev/pts        devpts  gid=5,mode=620 0 0\\
none            /dev/shm        tmpfs           defaults 0 0\\
none            /proc           proc            defaults 0 0\\
none            /sys            sysfs           efaults 0 0}}

\item {\bfseries umount the image file}\\
{\small {\tt \# umount /mnt/guest}}
\end{enumerate}

Now, the guest OS image {\small {\tt hd.img}} is ready. You can also reference {\small {\tt http://free.oszoo.org}} for quickstart images. But make sure to install the boot loader.

\subsection{Install Windows into an Image File using a VMX guest}
In order to install a Windows OS, you should keep {\small {\tt acpi=0}} in your VMX configuration file.

\section{VMX Guests}
\subsection{Editing the Xen VMX config file}
Make a copy of the example VMX configuration file {\small {\tt /etc/xen/xmeaxmple.vmx}} and edit the line that reads

{\small {\tt disk = [ 'file:/var/images/\emph{guest.img},ioemu:hda,w' ]}}

replacing \emph{guest.img} with the name of the guest OS image file you just made.

\subsection{Creating VMX guests}
Simply follow the usual method of creating the guest, using the -f parameter and providing the filename of your VMX configuration file:\\

{\small {\tt \# xend start\\
\# xm create /etc/xen/vmxguest.vmx}}

In the default configuration, VNC is on and SDL is off. Therefore VNC windows will open when VMX guests are created. If you want to use SDL to create VMX guests, set {\small {\tt sdl=1}} in your VMX configuration file. You can also turn off VNC by setting {\small {\tt vnc=0}}.
 
\subsection{Destroy VMX guests}
VMX guests can be destroyed in the same way as can paravirtualized guests. We recommend that you type the command 

{\small {\tt poweroff}} 

in the VMX guest's console first to prevent data loss. Then execute the command 

{\small {\tt xm destroy \emph{vmx\_guest\_id} }} 

at the Domain0 console.

\subsection{VMX window (X or VNC) Hot Key}
If you are running in the X environment after creating a VMX guest, an X window is created. There are several hot keys for control of the VMX guest that can be used in the window.
 
{\bfseries Ctrl+Alt+2} switches from guest VGA window to the control window. Typing {\small {\tt help }} shows the control commands help. For example, 'q' is the command to destroy the VMX guest.\\
{\bfseries Ctrl+Alt+1} switches back to VMX guest's VGA.\\
{\bfseries Ctrl+Alt+3} switches to serial port output. It captures serial output from the VMX guest. It works only if the VMX guest was configured to use the serial port. \\

\subsection{Save/Restore and Migration}
VMX guests currently cannot be saved and restored, nor migrated. These features are currently under active development.

\chapter{Vnets - Domain Virtual Networking}

Xen optionally supports virtual networking for domains using {\em vnets}.
These emulate private LANs that domains can use. Domains on the same
vnet can be hosted on the same machine or on separate machines, and the
vnets remain connected if domains are migrated. Ethernet traffic 
on a vnet is tunneled inside IP packets on the physical network. A vnet is a virtual
network and addressing within it need have no relation to addressing on 
the underlying physical network. Separate vnets, or vnets and the physical network,
can be connected using domains with more than one network interface and
enabling IP forwarding or bridging in the usual way.

Vnet support is included in \texttt{xm} and \xend:
\begin{verbatim}
# xm vnet-create <config>
\end{verbatim}
creates a vnet using the configuration in the file \verb|<config>|.
When a vnet is created its configuration is stored by \xend and the vnet persists until it is
deleted using
\begin{verbatim}
# xm vnet-delete <vnetid>
\end{verbatim}
The vnets \xend knows about are listed by
\begin{verbatim}
# xm vnet-list
\end{verbatim}
More vnet management commands are available using the
\texttt{vn} tool included in the vnet distribution.

The format of a vnet configuration file is
\begin{verbatim}
(vnet (id       <vnetid>)
      (bridge   <bridge>)
      (vnetif   <vnet interface>)
      (security <level>))
\end{verbatim}
White space is not significant. The parameters are:
\begin{itemize}
  \item \verb|<vnetid>|: vnet id, the 128-bit vnet identifier. This can be given
    as 8 4-digit hex numbers separated by colons, or in short form as a single 4-digit hex number.
    The short form is the same as the long form with the first 7 fields zero.
    Vnet ids must be non-zero and id 1 is reserved.

  \item \verb|<bridge>|: the name of a bridge interface to create for the vnet. Domains
    are connected to the vnet by connecting their virtual interfaces to the bridge.
    Bridge names are limited to 14 characters by the kernel.

  \item \verb|<vnetif>|: the name of the virtual interface onto the vnet (optional). The
    interface encapsulates and decapsulates vnet traffic for the network and is attached
    to the vnet bridge. Interface names are limited to 14 characters by the kernel.

  \item \verb|<level>|: security level for the vnet (optional). The level may be one of 
      \begin{itemize}
        \item \verb|none|: no security (default). Vnet traffic is in clear on the network.
        \item \verb|auth|: authentication. Vnet traffic is authenticated using IPSEC
           ESP with hmac96.
        \item \verb|conf|: confidentiality. Vnet traffic is authenticated and encrypted
           using IPSEC ESP with hmac96 and AES-128.
      \end{itemize}
      Authentication and confidentiality are experimental and use hard-wired keys at present.
\end{itemize}
When a vnet is created its configuration is stored by \xend and the vnet persists until it is
deleted using \texttt{xm vnet-delete <vnetid>}. The interfaces and bridges used by vnets
are visible in the output of \texttt{ifconfig} and \texttt{brctl show}.

\section{Example}
If the file \path{vnet97.sxp} contains
\begin{verbatim}
(vnet (id 97) (bridge vnet97) (vnetif vnif97)
      (security none))
\end{verbatim}
Then \texttt{xm vnet-create vnet97.sxp} will define a vnet with id 97 and no security.
The bridge for the vnet is called vnet97 and the virtual interface for it is vnif97.
To add an interface on a domain to this vnet set its bridge to vnet97
in its configuration. In Python:
\begin{verbatim}
vif="bridge=vnet97"
\end{verbatim}
In sxp:
\begin{verbatim}
(dev (vif (mac aa:00:00:01:02:03) (bridge vnet97)))
\end{verbatim}
Once the domain is started you should see its interface in the output of \texttt{brctl show}
under the ports for \texttt{vnet97}.

To get best performance it is a good idea to reduce the MTU of a domain's interface
onto a vnet to 1400. For example using \texttt{ifconfig eth0 mtu 1400} or putting
\texttt{MTU=1400} in \texttt{ifcfg-eth0}.
You may also have to change or remove cached config files for eth0 under
\texttt{/etc/sysconfig/networking}. Vnets work anyway, but performance can be reduced
by IP fragmentation caused by the vnet encapsulation exceeding the hardware MTU.

\section{Installing vnet support}
Vnets are implemented using a kernel module, which needs to be loaded before
they can be used. You can either do this manually before starting \xend, using the
command \texttt{vn insmod}, or configure \xend to use the \path{network-vnet}
script in the xend configuration file \texttt{/etc/xend/xend-config.sxp}:
\begin{verbatim}
(network-script        network-vnet)
\end{verbatim}
This script insmods the module and calls the \path{network-bridge} script.

The vnet code is not compiled and installed by default.
To compile the code and install on the current system
use \texttt{make install} in the root of the vnet source tree,
\path{tools/vnet}. It is also possible to install to an installation
directory using \texttt{make dist}. See the \path{Makefile} in
the source for details.

The vnet module creates vnet interfaces \texttt{vnif0002},
\texttt{vnif0003} and \texttt{vnif0004} by default. You can test that
vnets are working by configuring IP addresses on these interfaces
and trying to ping them across the network. For example, using machines
hostA and hostB:
\begin{verbatim}
hostA# ifconfig vnif0004 10.0.0.100 up
hostB# ifconfig vnif0004 10.0.0.101 up
hostB# ping 10.0.0.100
\end{verbatim}

The vnet implementation uses IP multicast to discover vnet interfaces, so
all machines hosting vnets must be reachable by multicast. Network switches
are often configured not to forward multicast packets, so this often
means that all machines using a vnet must be on the same LAN segment,
unless you configure vnet forwarding.

You can test multicast coverage by pinging the vnet multicast address:
\begin{verbatim}
# ping -b 224.10.0.1
\end{verbatim}
You should see replies from all machines with the vnet module running.
You can see if vnet packets are being sent or received by dumping traffic
on the vnet UDP port:
\begin{verbatim}
# tcpdump udp port 1798
\end{verbatim}

If multicast is not being forwaded between machines you can configure
multicast forwarding using vn. Suppose we have machines hostA on 10.10.0.100
and hostB on 10.11.0.100 and that multicast is not forwarded between them.
We use vn to configure each machine to forward to the other:
\begin{verbatim}
hostA# vn peer-add hostB
hostB# vn peer-add hostA
\end{verbatim}
Multicast forwarding needs to be used carefully - you must avoid creating forwarding
loops. Typically only one machine on a subnet needs to be configured to forward,
as it will forward multicasts received from other machines on the subnet.

%% Chapter Glossary of Terms moved to glossary.tex
\chapter{Glossary of Terms}

\begin{description}

\item[BVT] The BVT scheduler is used to give proportional fair shares
  of the CPU to domains.

\item[Domain] A domain is the execution context that contains a
  running {\bf virtual machine}.  The relationship between virtual
  machines and domains on Xen is similar to that between programs and
  processes in an operating system: a virtual machine is a persistent
  entity that resides on disk (somewhat like a program).  When it is
  loaded for execution, it runs in a domain.  Each domain has a {\bf
    domain ID}.

\item[Domain 0] The first domain to be started on a Xen machine.
  Domain 0 is responsible for managing the system.

\item[Domain ID] A unique identifier for a {\bf domain}, analogous to
  a process ID in an operating system.

\item[Full virtualization] An approach to virtualization which
  requires no modifications to the hosted operating system, providing
  the illusion of a complete system of real hardware devices.

\item[Hypervisor] An alternative term for {\bf VMM}, used because it
  means `beyond supervisor', since it is responsible for managing
  multiple `supervisor' kernels.

\item[Live migration] A technique for moving a running virtual machine
  to another physical host, without stopping it or the services
  running on it.

\item[Paravirtualization] An approach to virtualization which requires
  modifications to the operating system in order to run in a virtual
  machine.  Xen uses paravirtualization but preserves binary
  compatibility for user space applications.

\item[Shadow pagetables] A technique for hiding the layout of machine
  memory from a virtual machine's operating system.  Used in some {\bf
  VMMs} to provide the illusion of contiguous physical memory, in
  Xen this is used during {\bf live migration}.

\item[Virtual Block Device] Persistant storage available to a virtual
  machine, providing the abstraction of an actual block storage device.
  {\bf VBD}s may be actual block devices, filesystem images, or
  remote/network storage.

\item[Virtual Machine] The environment in which a hosted operating
  system runs, providing the abstraction of a dedicated machine.  A
  virtual machine may be identical to the underlying hardware (as in
  {\bf full virtualization}, or it may differ, as in {\bf
  paravirtualization}).

\item[VMM] Virtual Machine Monitor - the software that allows multiple
  virtual machines to be multiplexed on a single physical machine.

\item[Xen] Xen is a paravirtualizing virtual machine monitor,
  developed primarily by the Systems Research Group at the University
  of Cambridge Computer Laboratory.

\item[XenLinux] A name for the port of the Linux kernel that
  runs on Xen.

\end{description}


\end{document}


%% Other stuff without a home

%% Instructions Re Python API

%% Other Control Tasks using Python
%% ================================

%% A Python module 'Xc' is installed as part of the tools-install
%% process. This can be imported, and an 'xc object' instantiated, to
%% provide access to privileged command operations:

%% # import Xc
%% # xc = Xc.new()
%% # dir(xc)
%% # help(xc.domain_create)

%% In this way you can see that the class 'xc' contains useful
%% documentation for you to consult.

%% A further package of useful routines (xenctl) is also installed:

%% # import xenctl.utils
%% # help(xenctl.utils)

%% You can use these modules to write your own custom scripts or you
%% can customise the scripts supplied in the Xen distribution.



% Explain about AGP GART


%% If you're not intending to configure the new domain with an IP
%% address on your LAN, then you'll probably want to use NAT. The
%% 'xen_nat_enable' installs a few useful iptables rules into domain0
%% to enable NAT. [NB: We plan to support RSIP in future]



%% Installing the file systems from the CD
%% =======================================

%% If you haven't got an existing Linux installation onto which you
%% can just drop down the Xen and Xenlinux images, then the file
%% systems on the CD provide a quick way of doing an install. However,
%% you would be better off in the long run doing a proper install of
%% your preferred distro and installing Xen onto that, rather than
%% just doing the hack described below:

%% Choose one or two partitions, depending on whether you want a
%% separate /usr or not. Make file systems on it/them e.g.:
%% mkfs -t ext3 /dev/hda3
%% [or mkfs -t ext2 /dev/hda3 && tune2fs -j /dev/hda3 if using an old
%% version of mkfs]

%% Next, mount the file system(s) e.g.:
%%   mkdir /mnt/root && mount /dev/hda3 /mnt/root
%%   [mkdir /mnt/usr && mount /dev/hda4 /mnt/usr]
  
%% To install the root file system, simply untar /usr/XenDemoCD/root.tar.gz:
%%   cd /mnt/root && tar -zxpf /usr/XenDemoCD/root.tar.gz

%% You'll need to edit /mnt/root/etc/fstab to reflect your file system
%% configuration. Changing the password file (etc/shadow) is probably a
%% good idea too.

%% To install the usr file system, copy the file system from CD on
%% /usr, though leaving out the "XenDemoCD" and "boot" directories:
%%   cd /usr && cp -a X11R6 etc java libexec root src bin dict kerberos
%%    local sbin tmp doc include lib man share /mnt/usr

%% If you intend to boot off these file systems (i.e. use them for
%% domain 0), then you probably want to copy the /usr/boot
%% directory on the cd over the top of the current symlink to /boot
%% on your root filesystem (after deleting the current symlink)
%% i.e.:
%%   cd /mnt/root ; rm boot ; cp -a /usr/boot .
